\documentclass[10pt,a4paper]{amsart}
\usepackage[margin=19mm]{geometry}
\usepackage{amsmath,amssymb,mathtools}
\usepackage[scr=rsfs,cal=euler]{mathalfa}
\usepackage{xfrac}
\usepackage[unicode,hidelinks,bookmarks=false]{hyperref}
\newcommand{\ie}{i.e.\ }
\newcommand{\cf}{cf.\ }
\newcommand{\resp}{respectively\ }
\newcommand{\Subsection}{\subsection}
\providecommand{\nobreakdash}{\nobreak\hbox{-}}
\setlength{\emergencystretch}{2em}
\multlinegap0pt
\usepackage{bm}
\usepackage{amssymb}
\usepackage{enumerate}
\usepackage[shortlabels]{enumitem}
\usepackage{csquotes}
\newtheorem{lemma}{Lemma}
\newcommand{\abs}[1]{{\left|#1\right|}}
\newcommand{\bN}{\mathbb{N}}
\newcommand{\br}{{\mathbf{r}}}
\newcommand{\by}{{\mathbf{y}}}
\newcommand{\cJ}{\mathcal{J}}
\newcommand{\cO}{\mathcal{O}}
\newcommand{\cP}{\mathcal{P}}
\newcommand{\pairs}{{\rm B}}
\newcommand{\red}{\mathrm{red}}
\title{\textsc{Limits of descent-biased trees}}
\author{Victor Dubach, Paul Th\'evenin, Stephan Wagner}
\date{Source: arXiv:2609.15451v1 (CC BY 4.0)}
\begin{document}
\maketitle

\noindent\emph{Source and licence.} \textsc{Limits of descent-biased trees}. Version arXiv:2609.15451v1 (CC BY 4.0), distributed by arXiv under the Creative Commons Attribution 4.0 International licence, http://creativecommons.org/licenses/by/4.0/, as recorded in the arXiv licence metadata for this exact version and shown on the abstract page https://arxiv.org/abs/2609.15451v1. Full licence notice: \texttt{licenses/arxiv-2609.15451-CC-BY-4.0.txt}.\par\smallskip

\noindent\textit{Editorial scope.} Counted unit: the article's lemma tilde (source0.tex lines 1485--1510). The article's complete original proof of it is delivered with it (source0.tex lines 1593--1704, one \texttt{proof} environment of 112 lines, which the article places away from the statement). No mathematics is written, repaired or re-derived: the statement and the proof are the article's own lines, in the article's own notation, and no retained line is substituted or re-flowed.  Retained uncounted as the source-local context the proof needs: lines 1348--1358; lines 1384--1386.  Nothing was omitted from any retained span, so the package carries no omission record.  No retained line contains a citation, so the wrapper carries no bibliography and no bibliography entry.  No retained line contains a figure environment, an image-inclusion command or drawing code, so no image is delivered and the package declares no asset dependency.  Editorial formatting only, in addition: an AMS \texttt{amsart} preamble that restates the article's own declarations and macros used by the retained lines, namely the environments \texttt{lemma} on separate counters, together with the standard packages those lines need.  The retained environments print in this document as Lemma 1 in order of appearance, whereas the article prints the same items with the numbering of its own full document.  The complete native source of the article as distributed in the arXiv e-print for this exact version is bundled unchanged as \texttt{sources/210417/source0.tex}.
\begin{align}\label{eq: equadiff on Fell with everything}
    \left( e^{-A(x)} - qx \right)
    \partial_x F_\ell(x,\by)
    = &\sum_{s=0}^\ell \sum_{\cJ\in \cP_{1,s}^\red} \frac{1}{s!} 
    M_\cJ(\by)
    \prod_{1\le k\le s} F_{\abs{J_k}}\left(\big. x , \varphi_{J_k}(\by) \right) \nonumber\\
    &+ q \sum_{s=0}^\ell \sum_{\cJ\in \cP_{2,s}^\red} \frac{1}{s!} 
    M_\cJ(\by)
    F_{\abs{J_0}}\left(\big. x , \psi_{J_0}(\by) \right)
    \prod_{1\le k\le s} F_{\abs{J_k}}\left(\big. x , \varphi_{J_k}(\by) \right) \,.
\end{align}

\begin{equation}\label{eq: claim singularity fellr}
    f_{\ell,\br}(x) \sim C_{\ell,\br} (1-x/x_0)^{\frac12 - \ell - \frac12 \sum \br}
\end{equation}

\begin{lemma}\label{lem: asymptotics second term Rellr tilde}
    If $\ell=2$, then as $x\to x_0$,
    \begin{equation*}
        \partial_\by^\br R_\ell(x,\by) \vert_{\by=\bm{1}}
        = \cO\left( (1-x/x_0)^{\frac12-\ell-\frac12\sum\br} \right).
    \end{equation*}
    If $\ell\ge3$, then as $x\to x_0$,
    \begin{equation*}
        \partial_\by^\br R_\ell(x,\by) \vert_{\by=\bm{1}}
        \;\sim\;
        q\, C_{\ell,\br}^{(2)}\,
        (1-x/x_0)^{-\ell-\frac12\sum\br},
    \end{equation*}
    where
    \begin{equation*}
        C_{\ell,\br}^{(2)} 
        = \sum_{\substack{ J_0 \sqcup J_1 = [0,\ell] \\ 0\in J_0,\, \abs{J_0}<\ell,\, \abs{J_1}<\ell }}
        \sum_{\substack{ \bm{0}\preceq \bm{\alpha}\preceq \br \\ \alpha_{i,j}=r_{i,j} \text{ for all } i,j\in J_0 \\  \alpha_{i,j}=0 \text{ for all } i,j\in J_1 }}
        \binom{\br}{\bm{\alpha}}
        C_{\abs{J_0}, \bm\alpha^{J_0,\ell+1}}
        C_{\abs{J_1}, (\br-\bm\alpha)^{0,J_1}} \,.
    \end{equation*}
    Here, for $(J_0,J_1)$ as above and $\bm{\alpha} \in \bN_0^{\pairs_\ell}$, we let $\bm{\alpha}^{J_0,\ell+1}$ denote the element of $\bN_0^{\pairs_{\abs{J_0}}}$ defined by identifying $[0,\abs{J_0}]$ with $J_0 \sqcup \{\ell+1\}$ (by the unique order-preserving bijection) and setting $\alpha^{J_0,\ell+1}_{i,j} := \alpha_{i,j}$ and $\alpha^{J_0,\ell+1}_{i,\ell+1} := \sum_{h\in J_1} \alpha_{i,h}$ for all $i,j\in J_0$.
    Moreover, we let $\bm{\alpha}^{0,J_1}$ denote the element of $\bN_0^{\pairs_{\abs{J_1}}}$ defined by identifying $[0,\abs{J_1}]$ with $\{0\} \sqcup J_1$ and setting $\alpha^{0,J_1}_{i,j} = \alpha_{i,j}$ and $\alpha^{0,J_1}_{0,j} = \sum_{h\in J_0} \alpha_{h,j}$ for all $i,j\in J_1$.
    We set $C_{\ell,\br}^{(2)} := 0$ if $\ell=2$.
\end{lemma}

\begin{proof}[Proof of Lemma~\ref{lem: asymptotics second term Rellr tilde}]
Recall that $R_{\ell}$ is defined by \eqref{eq: equadiff on Fell with everything} after taking out the terms containing the function $F_\ell$:
\begin{align}
    R_\ell(x,\by)
    = &\sum_{s=0}^\ell \sum_{\substack{\cJ\in \cP_{1,s}^\red \\ \abs{J_k}<\ell \text{ for all } k\in[1,s]}} \frac{1}{s!} 
    M_\cJ(\by)
    \prod_{1\le k\le s} F_{\abs{J_k}}\left(\big. x , \varphi_{J_k}(\by) \right) 
    \label{eq: def Rell first line}\\
    &+ q \sum_{s=0}^\ell \sum_{\substack{\cJ\in \cP_{2,s}^\red \\ \abs{J_k}<\ell \text{ for all } k\in[0,s]}} \frac{1}{s!} 
    M_\cJ(\by)
    F_{\abs{J_0}}\left(\big. x , \psi_{J_0}(\by) \right)
    \prod_{1\le k\le s} F_{\abs{J_k}}\left(\big. x , \varphi_{J_k}(\by) \right) \,. 
    \label{eq: def Rell second line}
\end{align}
Let us analyze the asymptotics, as $x\to x_0$, of each term appearing in $\partial_\by^\br R_\ell(x,\by) \vert_{\by=\bm{1}}$, starting with the first line \eqref{eq: def Rell first line} of the formula.
\begin{itemize}
    \item Let $s=0$ and $\cJ$ be the unique partition in $\cP_{1,0}^\red$.
    The corresponding term is $\partial_\by^\br M_\cJ(\by) \vert_{\by=\bm{1}} = \cO(1)$.
    \item Let $s=1$ and $\cJ\in\cP_{1,1}^\red$ be such that $\abs{J_1}<\ell$.
    Then $M_\cJ(\bm{1}) F_{\abs{J_1}}(x,\bm{1}) = \cO\left( (1-x/x_0)^{\frac12-\abs{J_1}} \right)$ by the induction hypothesis \eqref{eq: claim singularity fellr}.
    Furthermore, each differentiation step changes this exponent by $0$ or $-\frac12$, so that
    \[
        \partial_\by^\br \left. \left[ M_\cJ(\by) F_{\abs{J_1}}\left( x,\varphi_{J_1}(\by) \right) \right] \right\vert_{\by=\bm{1}} 
        = \cO\left( (1-x/x_0)^{\frac12-(\ell-1)-\frac12\sum\br} \right) \,.
    \]
    \item Let $s\ge2$ and $\cJ\in\cP_{1,s}^\red$.
    Then $M_\cJ(\bm{1}) \prod_{1\le k\le s} F_{\abs{J_k}}(x,\bm{1}) = \cO\left( (1-x/x_0)^{\frac{s}{2}-\sum_{1\le k\le s}\abs{J_k}} \right)$ by the induction hypothesis \eqref{eq: claim singularity fellr}.
    Furthermore, each differentiation step changes this exponent by $0$ or $-\frac12$, so that
    \[
        \partial_\by^\br \left. \left[ M_\cJ(\by) \prod_{1\le k\le s} F_{\abs{J_k}}\left( x,\varphi_{J_k}(\by) \right) \right] \right\vert_{\by=\bm{1}} 
        = \cO\left( (1-x/x_0)^{\frac{s}{2}-\ell-\frac12\sum\br} \right) \,.
    \]
\end{itemize}
Each case above yields a $o\left( (1-x/x_0)^{\frac12-\ell-\frac12\sum\br} \right)$.
Now we turn to the second line \eqref{eq: def Rell second line}, which is handled similarly.
\begin{itemize}
    \item Let $s=0$ and $\cJ\in\cP_{2,0}^\red$ be such that $\abs{J_0}<\ell$.
    Then $M_\cJ(\bm{1}) F_{\abs{J_0}}(x,\bm{1}) = \cO\left( (1-x/x_0)^{\frac12-\abs{J_0}} \right)$ by the induction hypothesis \eqref{eq: claim singularity fellr}.
    Furthermore, each differentiation step changes this exponent by $0$ or $-\frac12$, so that
    \[
        \partial_\by^\br \left. \left[ M_\cJ(\by) F_{\abs{J_0}}\left( x,\psi_{J_0}(\by) \right) \right] \right\vert_{\by=\bm{1}} 
        = \cO\left( (1-x/x_0)^{\frac12-(\ell-1)-\frac12\sum\br} \right) \,.
    \]
    \item Let $s=1$ and $\cJ\in\cP_{2,1}^\red$ be such that $\abs{J_0}<\ell$ and $\abs{J_1}<\ell$.
    Then, by the induction hypothesis \eqref{eq: claim singularity fellr}, $M_\cJ(\bm{1}) F_{\abs{J_0}}(x,\bm{1}) F_{\abs{J_1}}(x,\bm{1}) = \cO\left( (1-x/x_0)^{1-\abs{J_0}-\abs{J_1}} \right)$.
    Furthermore, each differentiation step changes this exponent by $0$ or $-\frac12$, so that
    \[
        \partial_\by^\br \left. \left[ M_\cJ(\by) F_{\abs{J_0}}\left( x,\psi_{J_0}(\by) \right) F_{\abs{J_1}}\left( x,\varphi_{J_1}(\by) \right) \right] \right\vert_{\by=\bm{1}} 
        = \cO\left( (1-x/x_0)^{1-(\ell+1)-\frac12\sum\br} \right) \,.
    \]
    Let us analyze this case with more care.
    Specifically, we wish to identify which $\cJ\in\cP_{2,1}^\red$ yield an optimal exponent, that is, $-\ell$, before differentiating.
    First, assume $\ell=2$.
    Then, $\abs{J_0}+\abs{J_1} = \ell+1 = 3$ would imply either $\abs{J_0}=2$ or $\abs{J_1}=2$, which is forbidden in \eqref{eq: def Rell second line}.
    So, if $\ell=2$, we can improve to
    \[
        \partial_\by^\br \left. \left[ M_\cJ(\by) F_{\abs{J_0}}\left( x,\psi_{J_0}(\by) \right) F_{\abs{J_1}}\left( x,\varphi_{J_1}(\by) \right) \right] \right\vert_{\by=\bm{1}} 
        = \cO\left( (1-x/x_0)^{1-\ell-\frac12\sum\br} \right) \,.
    \]
    Next, assume $\ell\ge3$.
    Then, there exist partitions $(J_0,J_1)$ of $[0,\ell]$ with $0\in J_0$, $\abs{J_0}<\ell$ and $\abs{J_1}<\ell$.
    For such a partition, we have $M_\cJ(\bm{1}) F_{\abs{J_0}}(x,\bm{1}) F_{\abs{J_1}}(x,\bm{1}) = \Theta\left( (1-x/x_0)^{-\ell} \right)$ by the induction hypothesis \eqref{eq: claim singularity fellr}.
    Furthermore, to reach the optimal exponent $-\ell-\frac12\sum\br$ after differentiating, one needs to differentiate one of the two $F$ terms at each differentiation step.
    Recall that the family $\psi_{J_0}(\by)$ is defined in terms of $(y_{i,j})_{\{i,j\}\in \pairs_\ell^{(J_0)}}$ and the family $\varphi_{J_1}(\by)$ is defined in terms of $(y_{i,j})_{\{i,j\}\in \pairs_\ell^{(J_1)}}$.
    Recall also that $J_0 \sqcup J_1 = [0,\ell]$.
    Therefore, when differentiating $r_{i,j}$ times with respect to $y_{i,j}$, there are the following cases:
    \begin{itemize}
        \item If $\{i,j\}\in \pairs_\ell^{(J_0)} \cap \pairs_\ell^{(J_1)}$, that is, if $i\in J_0$ and $j\in J_1$ (or vice versa), then we get to choose which $F$ term to differentiate.
        Differentiating $F_{\abs{J_0}}( x,\psi_{J_0}(\by) )$, resp.\ $F_{\abs{J_1}}( x,\varphi_{J_1}(\by) )$, with respect to $y_{i,j}$ yields
        \[
            \left( \prod_{h\in J_1 \setminus \{j\}} y_{i,h} \right) \left[ \partial_{i,\ell+1} F_{\abs{J_0}} \right]( x,\psi_{J_0}(\by) ) \,,
            \quad\text{resp.}\quad
            \left( \prod_{h\in J_0 \setminus \{i\}} y_{h,j} \right) \left[ \partial_{0,j} F_{\abs{J_1}} \right]( x,\psi_{J_1}(\by) ) \,.
        \]
        Here, the first derivative is with respect to the argument of $F_{\abs{J_0}}$ indexed by $\{i,\ell+1\}$, where we identify $[0,\abs{J_0}]$ with $J_0 \sqcup \{\ell+1\}$.
        The second derivative is with respect to the argument of $F_{\abs{J_1}}$ indexed by $\{0,j\}$, where we identify $[0,\abs{J_1}]$ with $\{0\} \sqcup J_1$.
        \item Otherwise, $\{i,j\} \in \pairs_\ell^{(J_0)}$ if and only if both $i$ and $j$ are in $J_0$.
        In that case, we have to differentiate the $F_\abs{J_0}$ term.
        Differentiating $F_{\abs{J_0}}( x,\psi_{J_0}(\by) )$ with respect to $y_{i,j}$ yields
        \[
            \left[ \partial_{i,j} F_{\abs{J_0}} \right]( x,\psi_{J_0}(\by) ) \,,
        \]
        where the first derivative is with respect to the argument of $F_{\abs{J_0}}$ indexed by $\{i,j\}$, where we identify $[0,\abs{J_0}]$ with $J_0 \sqcup \{\ell+1\}$.
        \item Likewise, if both $i$ and $j$ are in $J_1$, we have to differentiate the $F_\abs{J_1}$ term.
        The function $F_{\abs{J_1}}( x,\varphi_{J_1}(\by) )$ has derivative with respect to $y_{i,j}$ given by
        \[
            \left[ \partial_{i,j} F_{\abs{J_1}} \right]( x,\varphi_{J_1}(\by) ) \,.
        \]
        Here, the derivative is with respect to the argument of $F_{\abs{J_1}}$ indexed by $\{i,j\}$, where we identify $[0,\abs{J_1}]$ with $\{0\} \sqcup J_1$.
    \end{itemize}
    In the end we get, with the notation of the lemma,
    \begin{multline*}
        \partial_\by^\br {\left. \left[ M_\cJ(\by) F_{\abs{J_0}}\left( x,\psi_{J_0}(\by) \right) F_{\abs{J_1}}\left( x,\varphi_{J_1}(\by) \right) \right] \right\vert_{\by=\bm{1}} } 
        \\\sim
        \sum_{\substack{ \bm{0}\preceq \bm{\alpha}\preceq \br \\ \alpha_{i,j}=r_{i,j} \text{ for all } i,j\in J_0 \\ \alpha_{i,j}=0 \text{ for all } i,j\in J_1}}
        %\left( \prod_{\{i,j\}\in \pairs_\ell} \binom{r_{i,j}}{\alpha_{i,j}} \right)
        \binom{\br}{\bm{\alpha}}
        f_{\abs{J_0}, \bm\alpha^{J_0,\ell+1}}(x)
        f_{\abs{J_1}, (\br-\bm\alpha)^{0,J_1}}(x) 
    \end{multline*}
    \item Let $s\ge2$ and $\cJ\in\cP_{2,s}^\red$.
    Then $M_\cJ(\bm{1}) F_{\abs{J_0}}(x,\bm{1}) \prod_{1\le k\le s} F_{\abs{J_k}}(x,\bm{1}) = \cO\left( (1-x/x_0)^{\frac{s+1}{2}-\sum_{1\le k\le s}\abs{J_k}} \right)$ by the induction hypothesis \eqref{eq: claim singularity fellr}.
    Furthermore, each differentiation step changes this exponent by $0$ or $-\frac12$, so that
    \[
        \partial_\by^\br \left. \left[ M_\cJ(\by) F_{\abs{J_0}}\left( x,\psi_{J_0}(\by) \right) \prod_{1\le k\le s} F_{\abs{J_k}}\left( x,\varphi_{J_k}(\by) \right) \right] \right\vert_{\by=\bm{1}} 
        = \cO\left( (1-x/x_0)^{\frac{s+1}{2}-(\ell+1)-\frac12\sum\br} \right) \,.
    \]
\end{itemize}
If $\ell=2$, then each case above yields a $\cO\left( (1-x/x_0)^{\frac12-\ell-\frac12\sum\br} \right)$.
If $\ell\ge3$, then the penultimate case yields a $\Theta\left( (1-x/x_0)^{-\ell-\frac12\sum\br} \right)$, while the other cases are negligible.
Using the induction hypothesis \eqref{eq: claim singularity fellr}, the conclusion follows.
\end{proof}

\end{document}
